\subsection{The bridge between accumulation and antiderivatives}

Suppose
\[
A(x)=\int_a^x f(t)\,dt.
\]
Increasing $x$ by a small amount $h$ adds approximately a thin contribution of width $h$ and height $f(x)$. Thus
\[
A(x+h)-A(x)\approx f(x)h,
\]
so the quotient
\[
\frac{A(x+h)-A(x)}{h}
\]
should approach $f(x)$. Accumulation therefore creates a function whose derivative recovers the local rate.

\begin{theorembox}[Fundamental theorem of calculus]
If $f$ is continuous and
\[
A(x)=\int_a^x f(t)\,dt,
\]
then $A'(x)=f(x)$. Consequently, if $F'(x)=f(x)$ on an interval, then
\[
\int_a^b f(x)\,dx=F(b)-F(a).
\]
\end{theorembox}

\begin{remarkbox}[Proof status]
The thin-strip argument explains the mechanism of the theorem. A complete proof requires estimates controlling the error in the approximation and is not developed in full here.
\end{remarkbox}

This is why the theorem is fundamental: a total defined by a limit of sums can be computed using a function whose derivative is the local contribution.

\begin{examplebox}[Area under a simple graph]
For $f(x)=x$ on $[0,2]$,
\[
\int_0^2 x\,dx=\left[\frac{x^2}{2}\right]_0^2=2.
\]
This agrees with the area of a triangle of base $2$ and height $2$.
\end{examplebox}

\begin{examplebox}[Signed accumulation]
\[
\int_{-1}^{1}x\,dx=0.
\]
The positive and negative contributions cancel. The integral is signed area, not total geometric area.
\end{examplebox}

For a non-negative integrand, the definite integral is non-negative. This simple sign check is often enough to detect an incorrect calculation.
